\documentclass[11pt,reqno]{amsart}

\usepackage[T1]{fontenc}
\usepackage{lmodern}
\usepackage{microtype}
\usepackage{mathtools}
\usepackage{amssymb}
\usepackage{booktabs}
\usepackage{enumitem}
\usepackage[hidelinks]{hyperref}
\usepackage[nameinlink,noabbrev]{cleveref}

\hypersetup{
  pdftitle={Clique partitions of chordal graphs with linear error},
  pdfauthor={Anonymous},
  pdfsubject={Erdos Problem 81},
  pdfkeywords={chordal graphs, clique partitions, fractional packings, stability}
}

\allowdisplaybreaks
\numberwithin{equation}{section}

\newtheorem{theorem}{Theorem}[section]
\newtheorem{proposition}[theorem]{Proposition}
\newtheorem{lemma}[theorem]{Lemma}
\newtheorem{corollary}[theorem]{Corollary}
\newtheorem{claim}[theorem]{Claim}
\theoremstyle{definition}
\newtheorem{definition}[theorem]{Definition}
\newtheorem{remark}[theorem]{Remark}

\newcommand{\cp}{\operatorname{cp}}
\newcommand{\cple}[1]{\operatorname{cp}_{\leq #1}}
\newcommand{\Wint}{W_4}
\newcommand{\Wfrac}{W_4^*}
\newcommand{\PhiF}{\Phi}
\newcommand{\distsp}{d_{\mathrm{sp}}}
\newcommand{\N}{\mathbb{N}}
\newcommand{\R}{\mathbb{R}}
\newcommand{\cS}{\mathcal{S}}
\newcommand{\cP}{\mathcal{P}}
\newcommand{\one}{\mathbf{1}}

\title[Linear-error clique partitions of chordal graphs]
  {Clique partitions of chordal graphs with linear error}

\author{Anonymous}
\date{September 7, 2026}

\subjclass[2020]{05C70, 05C35, 90C27}
\keywords{chordal graph, clique partition, fractional packing,
  symmetrization, stability}

\begin{document}

\begin{abstract}
We prove that the edges of every chordal graph on $n$ vertices can be
partitioned into $n^2/6+O(n)$ complete subgraphs.  This resolves a problem of
Erd\H{o}s, Ordman, and Zalcstein.  In fact, for all sufficiently large $n$, the
maximum clique-partition number among $n$-vertex chordal graphs is
\[
  \left\lfloor\frac{n(n+1)}6\right\rfloor .
\]
The proof combines a weighted fractional packing of triangles and copies of
$K_4$ with a quantitative stability argument.  Near the extremal
complete-split graphs we construct a partition using only edges and triangles,
with a strict saving for every structural defect.  Far from that family, a
finite-family packing transfer absorbs its subquadratic rounding error.  A
single-vertex symmetrization path and a first-entry argument connect the two
regimes.
\end{abstract}

\maketitle

\section{Introduction}

A \emph{clique partition} of a finite simple graph $G$ is a collection of
complete subgraphs whose edge sets partition $E(G)$.  We write $\cp(G)$ for
the minimum number of parts in such a partition.  Erd\H{o}s, Ordman, and
Zalcstein asked whether every chordal graph on $n$ vertices satisfies
\begin{equation}\label{eq:erdos-question}
  \cp(G)\leq \frac{n^2}{6}+O(n).
\end{equation}
Their complete-split examples show that the coefficient $1/6$ cannot be
improved \cite{ErdosOrdmanZalcstein1993}.  The question is listed as Erd\H{o}s
Problem 81 \cite{ErdosProblems81}.

Recall that a graph is chordal if it contains no induced cycle of length at
least four.  Split graphs, whose vertices can be divided into a clique and an
independent set, form an important subclass.  Chen, Erd\H{o}s, and Ordman
proved the bound $3n^2/16+O(n)$ for split graphs
\cite{ChenErdosOrdman1994}.  Two recent preprints give complementary partial
progress: Traverso proves \eqref{eq:erdos-question} for split graphs and
determines an associated fractional triangle functional, while Cipollini
proves $n^2/6+o(n^2)$ for all chordal graphs
\cite{Traverso2026,Cipollini2026}.  Neither partial result supplies the
linear-error conclusion for general chordal graphs.

Our main result is stronger than \eqref{eq:erdos-question}.

\begin{theorem}\label{thm:main}
There is an integer $N$ such that, for every $n\geq N$,
\[
  \max\{\cp(G): G\text{ is chordal and }|V(G)|=n\}
  =M(n):=\left\lfloor\frac{n(n+1)}6\right\rfloor .
\]
\end{theorem}

\begin{corollary}\label{cor:erdos81}
There is an absolute constant $C$ such that every chordal graph $G$ on $n$
vertices satisfies
\[
  \cp(G)\leq \frac{n^2}{6}+Cn.
\]
\end{corollary}

The proof has two distinct parts.  We first introduce a mixed packing
functional in which a triangle has gain $2$ and a copy of $K_4$ has gain $5$.
The corresponding fractional functional is well behaved under vertex
copying.  Repeated copies lead monotonically to a complete-split graph.  This
gives the correct global quadratic bound, but a general packing-transfer
theorem loses $o(n^2)$ and therefore cannot by itself prove
\eqref{eq:erdos-question}.

The second part repairs precisely that loss near the extremal family.  Starting
from a clique of order close to $n/3$, we regularize the choice of root clique
without changing the graph.  A balanced Vizing edge-colouring and the
H\"aggkvist--Janssen list-edge-colouring theorem then produce an integral
triangle packing with a fixed saving for both types of defect: edges outside
the root and missing edges between the root and its complement.  This yields a
local stability inequality in the actual labelled edit distance.  Finally, a
single-vertex version of the copying argument forces every graph with nearly
extremal fractional value into that local neighbourhood.  The first path
vertex entering the neighbourhood would otherwise satisfy incompatible upper
and lower bounds on its edit distance.

The constants used below are intentionally conservative.  Put
\begin{equation}\label{eq:constants}
  \rho=10^{-12},\qquad \eta=10^{-30}.
\end{equation}
For the packing-transfer theorem stated in \cref{thm:packing-transfer}, choose
a threshold $T(\varepsilon)$ that works uniformly for all graphs, and set
\begin{equation}\label{eq:N-definition}
  N=\max\{10^{32},T(\eta/2)\}.
\end{equation}
No numerical value for $T(\eta/2)$ is asserted.

\section{Preliminaries and external inputs}

All graphs in the paper are finite and simple.  For $r\geq2$, let
$\cple{r}(G)$ be the minimum number of parts in a clique partition of $G$ in
which every part has at most $r$ vertices.  Thus
\[
  \cp(G)\leq \cple{r}(G).
\]
We allow copies of $K_2$ as parts.  We write $e(G)=|E(G)|$,
$\Delta(G)$ for maximum degree, and $\omega(G)$ for clique number.

We use the following three published results.  Their quantifiers are recorded
explicitly because the uniformity in the third result is important.

\begin{theorem}[Vizing \cite{Vizing1964}]\label{thm:vizing}
Every finite simple graph $H$ has a proper edge-colouring with at most
$\Delta(H)+1$ colours.
\end{theorem}

\begin{theorem}[H\"aggkvist--Janssen \cite{HaggkvistJanssen1997}]
\label{thm:HJ}
The list-chromatic index of $K_p$ is at most $p$.  Equivalently, if every edge
of $K_p$ is assigned a list of at least $p$ colours, then the edges admit a
proper colouring in which each edge receives a colour from its own list.
\end{theorem}

For a graph $G$, consider edge-disjoint packings by triangles and copies of
$K_4$.  Assign gain $2$ to a triangle and gain $5$ to a copy of $K_4$.  Let
$\Wint(G)$ be the maximum total gain of an integral packing and let
$\Wfrac(G)$ be the optimum of the fractional relaxation.  Explicitly, if
$\mathcal{T}(G)$ and $\mathcal{K}_4(G)$ denote the triangles and copies of
$K_4$ in $G$, then
\begin{equation}\label{eq:mixed-primal}
 \begin{aligned}
 \Wfrac(G)=\max\quad
   &2\sum_{T\in\mathcal{T}(G)}x_T
    +5\sum_{K\in\mathcal{K}_4(G)}z_K,\\
 \text{subject to}\quad
   &x_T,z_K\geq0,\\
   &\sum_{T\ni e}x_T+\sum_{K\ni e}z_K\leq1
       &&\text{for every }e\in E(G).
 \end{aligned}
\end{equation}

\begin{theorem}[finite-family packing transfer]
\label{thm:packing-transfer}
For every $\varepsilon>0$ there is an integer $T(\varepsilon)$ such that every
graph $G$ with $|V(G)|\geq T(\varepsilon)$ satisfies
\[
  0\leq \Wfrac(G)-\Wint(G)\leq
  \varepsilon |V(G)|^2.
\]
The threshold is uniform over all graphs of the given order.
\end{theorem}

This is the specialization of the weighted finite-family theorem
\cite[Theorem~3.4]{RohatgiUrschelWellens2021} to $\{K_3,K_4\}$ and gains
$2,5$.  The proof of
\cite[Theorem~3.6]{RohatgiUrschelWellens2021} also makes the required uniform
application explicit.

We shall use standard properties of chordal graphs: induced subgraphs are
chordal; chordal graphs admit perfect elimination orderings; adding a vertex
whose neighbourhood is a clique preserves chordality; and every noncomplete
chordal graph has two nonadjacent simplicial vertices
\cite{Dirac1961,BlairPeyton1993}.  We also need the following convenient form
of the elimination-ordering property.

\begin{lemma}\label{lem:peo-ending-clique}
If $P$ is a clique in a chordal graph $G$, then $G$ has a perfect elimination
ordering whose final vertices are precisely the vertices of $P$.
\end{lemma}

\begin{proof}
As long as a vertex remains outside $P$, remove a simplicial vertex outside
$P$.  Such a vertex exists: if the current graph is complete, every outside
vertex is simplicial; otherwise, two nonadjacent simplicial vertices exist and
they cannot both lie in the clique $P$.  The induced graph left after each
removal is chordal.  Appending the vertices of $P$ in any order completes the
desired ordering.
\end{proof}

The fractional mixed packing has the dual formulation
\begin{equation}\label{eq:mixed-dual}
  \Wfrac(G)=\min\left\{
  \sum_{e\in E(G)}y_e:
  \begin{array}{l}
    y_e\geq0,\\[2pt]
    \sum_{e\in E(T)}y_e\geq2\quad\text{for every triangle }T,\\[2pt]
    \sum_{e\in E(K)}y_e\geq5\quad\text{for every }K\cong K_4
  \end{array}\right\}.
\end{equation}
Both primal and dual optima are attained.  For example, a dual coordinate
larger than $5$ can be replaced by $5$, so the dual minimum may be taken over
a compact box.

Define
\begin{equation}\label{eq:functional-definitions}
  \PhiF(G)=e(G)-\Wfrac(G),\qquad
  Q(n)=\frac{(2n+1)^2}{24}.
\end{equation}
Since the gain of a clique is the number of edges saved when it replaces its
constituent $K_2$ parts, we have the exact identity
\begin{equation}\label{eq:cp4-identity}
  \cple{4}(G)=e(G)-\Wint(G)
  =\PhiF(G)+\Wfrac(G)-\Wint(G).
\end{equation}
Moreover,
\begin{equation}\label{eq:Phi-cp3}
  \PhiF(G)\leq \cple{3}(G),
\end{equation}
because every integral triangle packing is an admissible mixed fractional
packing.  Notice that \eqref{eq:Phi-cp3} does not assert
$\PhiF(G)\leq\cp(G)$.

Finally,
\begin{equation}\label{eq:floor-identity}
  \lfloor Q(n)\rfloor
  =\left\lfloor\frac{n(n+1)}6\right\rfloor=M(n).
\end{equation}
Indeed, $Q(n)=n(n+1)/6+1/24$, while the fractional part of $n(n+1)/6$ is
either $0$ or $1/3$.

\section{An integral construction around a root clique}

Let $P$ be a clique in a chordal graph $G$, and introduce the data
\begin{equation}\label{eq:root-data}
  p=|P|,\quad H=G-P,\quad q=|V(H)|,\quad
  m=e(H),\quad w=\omega(H).
\end{equation}
Write
\begin{equation}\label{eq:root-defects}
  A=pq-e_G(P,V(H)),\qquad
  D=\max_{x\in P}|V(H)\setminus N_G(x)|.
\end{equation}
Thus $A$ is the total number of missing root--outside incidences and $D$ is
the largest missing column; here $e_G(X,Y)$ denotes the number of edges with
one endpoint in each of the disjoint sets $X$ and $Y$.

\begin{lemma}[terminal construction]\label{lem:terminal-construction}
Assume $p\geq1$, and put
\[
  c=\max\{p,\Delta(H)+1\},\qquad
  t=\left\lceil\frac{m}{c}\right\rceil.
\]
If
\begin{equation}\label{eq:host-condition}
  q-2D-4t\geq p,
\end{equation}
then
\begin{equation}\label{eq:terminal-bound}
  \cple{3}(G)
  \leq pq-\binom p2
    +\left(1-\frac{2p}{c}\right)m
    -\left(1-\frac{2(w-1)}p\right)A.
\end{equation}
The clique $P$ need not be maximum.
\end{lemma}

\begin{proof}
By \cref{thm:vizing}, the edges of $H$ have a proper colouring with $c$
colours.  Choose such a colouring for which the sum of the squares of the
colour-class sizes is minimum.  Its classes differ in size by at most one.
Indeed, if two classes differed by at least two, their union would contain an
alternating path with one more edge of the larger colour; interchanging the
two colours on that path would decrease the sum of squares.  Consequently
each class has at most $t$ edges.

Retain the $p$ largest classes.  Together they contain at least $pm/c$ edges.
We assign these classes bijectively to the vertices of $P$.  An outside edge
$uv$ assigned to $x\in P$ is retained when $x$ is adjacent to both $u$ and
$v$; it then yields the triangle $uvx$.

Choose a perfect elimination ordering ending in $P$, as in
\cref{lem:peo-ending-clique}, and orient every edge of $H$ from its earlier to
its later endpoint.  If $uv$ is oriented from $u$ to $v$, then
\[
  N_G(u)\cap P\subseteq N_G(v)\cap P,
\]
because the later neighbours of $u$ form a clique.  Hence the number of root
vertices that are invalid hosts for $uv$ is exactly
$a_u=p-|N_G(u)\cap P|$.  The vertex $u$ has at most $w-1$ later neighbours in
$H$, and therefore the total number of invalid-host incidences over all
outside edges is at most
\[
  (w-1)\sum_{u\in V(H)}a_u=(w-1)A.
\]
This remains an upper bound after only the selected colour classes are kept.

Choose the bijection from the selected classes to $P$ uniformly at random.
Each invalid host is assigned to its class with probability $1/p$, so some
bijection retains at least
\begin{equation}\label{eq:f-lower}
  f\geq \frac{pm}{c}-\frac{(w-1)A}{p}
\end{equation}
outside edges.  Use the corresponding $f$ triangles.  Since every selected
class is a matching of size at most $t$, each root vertex is incident with at
most $2t$ outside spokes already used by these triangles.

It remains to cover the edges within $P$.  For each $xy\in E(G[P])$, let
$L(xy)$ be the set of outside vertices $z$ such that $xyz$ is a triangle and
neither $xz$ nor $yz$ has already been used.  The union bound gives
\[
  |L(xy)|\geq q-2D-4t\geq p.
\]
Apply \cref{thm:HJ} to the complete graph on $P$ with these actual lists.  Its
proper list-edge-colouring assigns a host $z\in L(xy)$ to every root edge
$xy$.  The resulting root triangles are pairwise edge-disjoint and are also
edge-disjoint from the $f$ outside-edge triangles.

Use $K_2$ parts for every crossing or outside edge not yet covered.  The exact
number of parts is
\[
  pq-\binom p2+m-A-2f.
\]
Substituting \eqref{eq:f-lower} proves \eqref{eq:terminal-bound}.
\end{proof}

We shall repeatedly use the following immediate consequence.

\begin{corollary}\label{cor:strict-terminal-bound}
Under the hypotheses of \cref{lem:terminal-construction}, if in addition
\[
  c\leq\frac{9p}{5}
  \quad\text{and}\quad
  \frac{2(w-1)}p\leq\frac12,
\]
then
\begin{equation}\label{eq:strict-terminal-bound}
  \cple{3}(G)\leq pq-\binom p2-\frac m9-\frac A2.
\end{equation}
\end{corollary}

\begin{proof}
In \eqref{eq:terminal-bound}, the coefficient of $m$ is at most $-1/9$ and
the coefficient of $A$ is at most $-1/2$.
\end{proof}

\section{Strict regularization of the root}

The construction in \cref{lem:terminal-construction} becomes useful only when
both coefficients in \eqref{eq:terminal-bound} are bounded away from zero.
The next elementary observation controls vertices that are not complete to the
root.

\begin{lemma}\label{lem:nonadjacency}
With the notation of \eqref{eq:root-data}--\eqref{eq:root-defects}, suppose
$u\in V(H)$ is nonadjacent to some $x\in P$.  Then
\[
  d_H(u)\leq D+w-2.
\]
\end{lemma}

\begin{proof}
The common neighbours of two nonadjacent vertices in a chordal graph form a
clique: two nonadjacent common neighbours would induce a $4$-cycle.  Thus
$N_H(u)\cap N_H(x)$, together with $u$, is a clique in $H$, and the
intersection has at most $w-1$ vertices.  Every remaining neighbour of $u$ in
$H$ lies in $V(H)\setminus N_G(x)$.  This set has at most $D$ vertices and
contains $u$ itself, so it contains at most $D-1$ neighbours of $u$.  Adding
the two estimates proves the claim.
\end{proof}

\begin{lemma}[strict root regularization]\label{lem:root-regularization}
Suppose that a chordal graph $G$ contains a clique $P$ whose data satisfy
\begin{equation}\label{eq:regularization-hypotheses}
  p\geq128,\qquad |q-2p|\leq\frac p{64},\qquad
  m+A\leq\frac{p^2}{65536}.
\end{equation}
Then $G$ contains a clique $P'$---not necessarily a maximum clique---such that,
for the corresponding data $p',q',m',A'$,
\begin{equation}\label{eq:regularized-bound}
  \cple{3}(G)
  \leq p'q'-\binom{p'}2-\frac{m'}9-\frac{A'}2.
\end{equation}
\end{lemma}

\begin{proof}
First demote every root vertex whose original missing column has more than
$p/64$ entries.  Let $R$ be the set of demoted vertices, let $r=|R|$, and put
$P_0=P\setminus R$.  This changes only the designation of the root: no edge is
edited and no clique-partition part is chosen.  From
\eqref{eq:regularization-hypotheses},
\begin{equation}\label{eq:demotion-size}
  r\leq\frac{64A}{p}\leq\frac p{1024},\qquad
  p_0=p-r\geq\frac{99p}{100},\qquad q_0=q+r,
\end{equation}
and the largest missing column of the retained root satisfies $D_0\leq p/64$.

Let $H_0=G-P_0$.  Since $R$ was contained in the original clique, its outside
edge count is exactly
\[
  m_0=m+e_G(R,V(H))+\binom r2.
\]
Using $q\leq129p/64$ and \eqref{eq:demotion-size}, we obtain
\begin{align}
  m_0
  &\leq \left(\frac1{65536}
      +\frac{129}{64\cdot1024}
      +\frac1{2\cdot1024^2}\right)p^2  \notag\\
  &=\frac{4161}{2097152}p^2
   <\frac{p^2}{400}.
  \label{eq:m0-bound}
\end{align}
If $w_0=\omega(H_0)$, then
$\binom{w_0}{2}\leq m_0$, and hence
\begin{equation}\label{eq:w0-bound}
  w_0-1\leq\sqrt{2m_0}<\frac p{10}.
\end{equation}

Define the fixed defect bound
\begin{equation}\label{eq:Dbar-definition}
  \overline D=\max\left\{D_0,q_0-\frac{7p_0}{4}\right\}.
\end{equation}
Equations \eqref{eq:regularization-hypotheses} and
\eqref{eq:demotion-size} give
\begin{equation}\label{eq:Dbar-bound}
  \overline D
  \leq\max\left\{\frac p{64},
        \left(\frac{17}{64}+\frac{11}{4096}\right)p\right\}
  =\frac{1099}{4096}p<\frac p3.
\end{equation}

Starting from $P_0$, promote an outside vertex whenever its degree in the
current outside graph is at least $7/4$ times the current root order.  Every
promotion is legal.  Indeed, by \cref{lem:nonadjacency}, a vertex not complete
to the current root has outside degree at most
\[
  \overline D+w_0-2<\frac{13p}{30},
\]
far below $7p_0/4$.  Thus every promoted vertex is complete to the current
root, which consequently remains a clique.

The bound $\overline D$ remains valid throughout the procedure.  If $u$ is
promoted when the outside and root orders are $q_i$ and $p_i$, respectively,
then its missing column after promotion has size at most
\[
  q_i-1-d_{H_i}(u)
  \leq q_i-\frac{7p_i}{4}-1
  \leq q_0-\frac{7p_0}{4}\leq\overline D.
\]
The old missing columns can only shrink, and the outside clique number cannot
increase.

The edges removed from the outside graph in different promotions are
disjointly accounted for.  If $h$ vertices are promoted, then
\begin{equation}\label{eq:promotion-count}
  h\leq\frac{4m_0}{7p_0}\leq\frac p{693}.
\end{equation}
Let $P'$ and $H'$ be the root and outside graph at termination.  Then
$\Delta(H')<7p'/4$, and therefore
\begin{equation}\label{eq:cprime-bound}
  c':=\max\{p',\Delta(H')+1\}
  \leq\left\lceil\frac{7p'}4\right\rceil
  \leq\frac{9p'}5.
\end{equation}
The last inequality holds for every integer $p'\geq15$, whereas here
$p'\geq p_0\geq127$.  Also, by \eqref{eq:w0-bound},
\begin{equation}\label{eq:wprime-bound}
  \frac{2(w'-1)}{p'}<\frac{20}{99}<\frac12.
\end{equation}

Put $t'=\lceil m'/c'\rceil$.  Since $m'\leq m_0$ and $c'\geq p_0$,
\begin{equation}\label{eq:tprime-bound}
  t'\leq\left\lceil\frac{m_0}{p_0}\right\rceil
  \leq\frac p{396}+1.
\end{equation}
Finally, each promotion decreases $q-p$ by two, so
\begin{align}
 q'-p'-2D'-4t'
 &\geq
 \left(\frac{63}{64}-\frac2{693}-\frac23-\frac1{99}\right)p-4 \notag\\
 &=\frac{4505}{14784}p-4
  >\frac p4-4\geq0.
 \label{eq:final-host-margin}
\end{align}
Thus the host condition \eqref{eq:host-condition} holds for $P'$.  Apply
\cref{cor:strict-terminal-bound} using \eqref{eq:cprime-bound} and
\eqref{eq:wprime-bound}; this gives \eqref{eq:regularized-bound}.
\end{proof}

\begin{remark}
The promotion threshold $7p/4$ is what leaves a fixed negative coefficient on
$m'$ in \eqref{eq:regularized-bound}.  A threshold of $2p$ would permit that
coefficient to vanish.  The exact formula for $m_0$ also shows that every edge
created in the outside graph by demoting a root vertex has been charged.
\end{remark}

\section{Quantitative local stability}

Fix $n$, put $a=\lfloor n/3\rfloor$, and, for each $A_0\subseteq V(G)$ with
$|A_0|=a$, let $S_{A_0}$ be the complete-split graph on $V(G)$ with clique
side $A_0$.  Thus two vertices are adjacent in $S_{A_0}$ if and only if at
least one of them lies in $A_0$.  Define the labelled edit distance
\begin{equation}\label{eq:edit-distance}
  \distsp(G)=\min_{|A_0|=a}
    |E(G)\mathbin{\triangle}E(S_{A_0})|.
\end{equation}
A single operation that replaces the neighbourhood of one vertex by the
neighbourhood of a nonadjacent vertex changes at most $n-2$ unordered pairs.
Consequently, $\distsp$ changes by at most $n-2$ in one such operation.

\begin{lemma}\label{lem:local-root}
Let $G$ be chordal of order $n\geq1000$.  If
\[
  \distsp(G)\leq\rho n^2,
\]
then $G$ has a clique satisfying the hypotheses of
\cref{lem:root-regularization}.
\end{lemma}

\begin{proof}
Choose $A_0$ in \eqref{eq:edit-distance} with edit count
$E\leq\rho n^2$.  Let $P$ be a maximum clique of the chordal graph $G[A_0]$,
write $p=|P|$, and put $u=a-p$.

A chordal graph on $a$ vertices with clique number $p$ has at most
\begin{equation}\label{eq:chordal-edge-bound}
  (p-1)a-\binom p2
\end{equation}
edges.  To see this, sum the number of later neighbours in a perfect
elimination ordering, bounding the successive terms by
$\min\{p-1,a-i\}$.  The number of pairs missing from $G[A_0]$ is therefore at
least $\binom{u+1}{2}$.  These pairs are counted by $E$, and hence
\begin{equation}\label{eq:u-bound}
  u\leq\sqrt{2E}\leq\sqrt{2\rho}\,n.
\end{equation}

Since $p=a-u$ and $q=n-p$, for $n\geq1000$ we have
\begin{align}
  p&\geq \frac n3-1-\sqrt{2\rho}\,n
      \geq\frac{33n}{100}\geq128,\label{eq:p-local-range}\\
  p&\leq\frac n3,\qquad
  0\leq q-2p=n-3p
      \leq3+3\sqrt{2\rho}\,n\leq\frac p{64}.
  \label{eq:q-local-range}
\end{align}

The edges in the old independent side and the nonedges between the two old
sides contribute at most $E$ defects altogether.  Every newly relevant pair
involving $A_0\setminus P$ contributes at most $un$.  Thus the outside-edge
count and missing-incidence count for the root $P$ satisfy
\begin{align}
  m+A
  &\leq E+un
   \leq(\rho+\sqrt{2\rho})n^2 \notag\\
  &\leq\left(10^{-12}+\frac{3}{2\cdot10^6}\right)n^2
   <\frac{(33/100)^2n^2}{65536}
   \leq\frac{p^2}{65536}.
\end{align}
Together with \eqref{eq:p-local-range}--\eqref{eq:q-local-range}, this is
exactly \eqref{eq:regularization-hypotheses}.
\end{proof}

\begin{lemma}[local stability]\label{lem:local-stability}
Under the hypotheses of \cref{lem:local-root}, put
\[
  \mathcal{D}=Q(n)-\PhiF(G).
\]
Then $\mathcal{D}\geq0$ and
\begin{equation}\label{eq:local-stability}
  \distsp(G)
  \leq9\mathcal{D}+n\sqrt{\frac{2\mathcal{D}}3}+n.
\end{equation}
Moreover,
\begin{equation}\label{eq:local-integral-conclusion}
  \cple{3}(G)\leq M(n).
\end{equation}
\end{lemma}

\begin{proof}
Apply \cref{lem:root-regularization}, and use primes for the resulting root
data.  By \eqref{eq:Phi-cp3} and \eqref{eq:regularized-bound},
\begin{equation}\label{eq:deficit-lower}
  \mathcal{D}
  \geq Q(n)-\left(p'q'-\binom{p'}2\right)
      +\frac{m'}9+\frac{A'}2.
\end{equation}
The first difference on the right has the exact form
\begin{equation}\label{eq:square-identity}
  Q(n)-\left(p'(n-p')-\binom{p'}2\right)
  =\frac{(6p'-2n-1)^2}{24}.
\end{equation}
It follows that
\begin{equation}\label{eq:defect-and-root-control}
  m'+A'\leq9\mathcal{D},
  \qquad
  \left|p'-\frac{2n+1}{6}\right|
  \leq\sqrt{\frac{2\mathcal{D}}3}.
\end{equation}
In particular, $\mathcal{D}\geq0$.

The graph $G$ differs from the complete-split graph with clique side $P'$ on
exactly $m'+A'$ pairs.  Changing that clique side to have order
$\lfloor n/3\rfloor$ requires at most
$|p'-\lfloor n/3\rfloor|$ vertex-role changes, each costing at most $n$ edge
edits.  Since
\[
  \left|\frac{2n+1}{6}-\left\lfloor\frac n3\right\rfloor\right|<1,
\]
equation \eqref{eq:local-stability} follows from
\eqref{eq:defect-and-root-control}.

Finally, \eqref{eq:regularized-bound} and \eqref{eq:square-identity} give
$\cple{3}(G)\leq Q(n)$.  The left side is an integer, so
\eqref{eq:floor-identity} yields \eqref{eq:local-integral-conclusion}.
\end{proof}

\section{A monotone path of single-vertex copies}

Let $u$ and $v$ be nonadjacent vertices of a graph $G$.  Write
$G_{u\to v}$ for the graph obtained by deleting the edges incident with $u$
and then making $N_{G_{u\to v}}(u)=N_G(v)$.  In particular, $u$ and $v$
remain nonadjacent.

\begin{lemma}[copy inequality]\label{lem:copy-inequality}
For every graph $G$ and every nonadjacent pair $u,v$,
\begin{equation}\label{eq:copy-inequality}
  \PhiF(G_{u\to v})+\PhiF(G_{v\to u})\geq2\PhiF(G).
\end{equation}
\end{lemma}

\begin{proof}
Choose an optimal dual cover $y$ in \eqref{eq:mixed-dual}.  To obtain a cover
of $G_{u\to v}$, keep every unaffected edge weight and give each new edge
$ux$ the old weight $y_{vx}$.  Every triangle or $K_4$ containing $u$ is then
a copy of one containing $v$ in $G$, while no clique contains both $u$ and
$v$.  Thus all dual constraints remain valid.

Let $s_u$ and $s_v$ be the sums of the $y$-weights on edges incident with $u$
and $v$.  The copied covers of $G_{u\to v}$ and $G_{v\to u}$ have respective
weights
\[
  \Wfrac(G)-s_u+s_v
  \quad\text{and}\quad
  \Wfrac(G)-s_v+s_u.
\]
They therefore give
\[
  \Wfrac(G_{u\to v})+\Wfrac(G_{v\to u})
  \leq2\Wfrac(G).
\]
On the other hand, the two copied edge counts sum to $2e(G)$.  Subtracting the
preceding inequality from this identity gives
\eqref{eq:copy-inequality}.
\end{proof}

Vertices with the same open neighbourhood are called \emph{false twins}.
Every false-twin class is independent.  Suppose $A_1$ and $A_2$ are distinct
false-twin classes consisting of simplicial vertices, with no edge between
the two classes.  Keep the graph outside $A_1\cup A_2$ fixed.  If
$r=|A_1|+|A_2|$, let $H_j$ be the graph in which $j$ of these $r$ vertices
have neighbourhood $N_G(A_1)$ and the other $r-j$ have neighbourhood
$N_G(A_2)$.  Both neighbourhoods are cliques, so every $H_j$ is chordal.
For $0<j<r$, \cref{lem:copy-inequality} gives
\begin{equation}\label{eq:discrete-convexity}
  \PhiF(H_{j-1})+\PhiF(H_{j+1})\geq2\PhiF(H_j).
\end{equation}
Thus $j\mapsto\PhiF(H_j)$ is discretely convex.  At any interior index, at
least one adjacent value is no smaller.  Once such a direction is chosen,
the successive differences in that direction remain nonnegative by
\eqref{eq:discrete-convexity}.  We may therefore move one vertex at a time to
the corresponding endpoint without ever decreasing $\PhiF$.

At the endpoint the two selected false-twin classes have merged.  No class
outside their union splits: all of its vertices have the same adjacency to a
representative of the target class and consequently acquire the same
adjacency to every copied vertex.  The number of false-twin classes has
strictly decreased.

We next identify the graphs at which this procedure stops.

\begin{lemma}[terminal characterization]\label{lem:terminal-characterization}
Let $G$ be chordal.  If every two nonadjacent simplicial vertices of $G$ have
the same open neighbourhood, then $G$ is complete-split; that is,
\[
  G\cong K_k\vee\overline K_{n-k}
\]
for some $0\leq k\leq n$.
\end{lemma}

\begin{proof}
Suppose first that $G$ is disconnected and has an edge.  Choose a nonisolated
simplicial vertex in a nontrivial component and a simplicial vertex in another
component.  Their neighbourhoods are disjoint and at least one is nonempty,
contrary to the hypothesis.  Hence a disconnected terminal graph is edgeless.
Complete graphs are also of the required form.

It remains to consider a connected, noncomplete graph.  Take a clique tree of
$G$ with at least two maximal-clique bags.  A vertex private to a leaf bag is
simplicial and occurs in no other bag.  Private vertices chosen from distinct
leaves are nonadjacent, and hence they have a common open neighbourhood, say
$K$.

Each leaf has exactly one private vertex.  Otherwise, the neighbourhood of a
private vertex in that leaf would contain another private vertex from the
same leaf, whereas the neighbourhood of a private vertex in a different leaf
would not.  It follows that every leaf bag is $K$ together with its unique
private vertex.

Every member of $K$ occurs in every leaf bag.  By the running-intersection
property, it then occurs in every bag, because every node of a finite tree lies
on a path between two leaves.  Thus $K$ is complete and universal to
$V(G)\setminus K$.  Each leaf-private vertex is isolated in $G-K$.

If $G-K$ had a nontrivial component, choose a nonisolated simplicial vertex of
that chordal component.  Its neighbourhood in $G$ is a clique, so it is also
simplicial in $G$.  It is nonadjacent to a leaf-private vertex, but its
neighbourhood properly contains $K$, again a contradiction.  Therefore $G-K$
is independent, as required.
\end{proof}

\begin{proposition}[fine monotone symmetrization]
\label{prop:fine-symmetrization}
Every chordal graph $G$ of order $n$ admits a sequence
\[
  G=G_0,G_1,\ldots,G_L=S_{k,n-k}
\]
of chordal graphs such that
\[
  \PhiF(G_{i+1})\geq\PhiF(G_i)
  \quad\text{and}\quad
  |E(G_{i+1})\mathbin{\triangle}E(G_i)|\leq n-2
\]
for every $i$.  One may take $L\leq n(n-1)$.
\end{proposition}

\begin{proof}
If the condition in \cref{lem:terminal-characterization} fails, choose two
nonadjacent simplicial vertices with different neighbourhoods and take their
false-twin classes.  The preceding discrete-convexity argument merges those
classes by a sequence of nondecreasing single-vertex copies.  Each copied
vertex changes at most $n-2$ incident pairs.  Repeat.  A full-class merge
strictly decreases the number of false-twin classes, so there are at most
$n-1$ merges, each requiring at most $n$ copies.  The terminal graph is
complete-split by \cref{lem:terminal-characterization}.
\end{proof}

The single-vertex resolution in \cref{prop:fine-symmetrization} is essential:
the proof below controls the first point at which the path crosses the boundary
of a fixed edit-distance neighbourhood.

\section{The mixed functional on complete-split graphs}

Write
\[
  S_{k,\ell}=K_k\vee\overline K_\ell,
  \qquad k+\ell=n.
\]

\begin{lemma}\label{lem:split-functional}
If $k\geq4$ and $\ell\geq1$, then
\begin{equation}\label{eq:split-functional}
  \PhiF(S_{k,\ell})=
  \max\left\{
    k\ell-\binom k2,
    \frac{2k\ell-\binom k2}{3},
    \frac{k\ell+\binom k2}{6}
  \right\}.
\end{equation}
\end{lemma}

\begin{proof}
Average an optimal dual cover over all automorphisms of $S_{k,\ell}$.  We may
therefore assign a common weight $x$ to each edge in the clique side and a
common weight $y$ to each spoke.  The triangle and $K_4$ constraints in
\eqref{eq:mixed-dual} reduce to
\begin{equation}\label{eq:split-dual-region}
  x\geq\frac56,\qquad y\geq0,\qquad
  x+2y\geq2,\qquad x+y\geq\frac53.
\end{equation}
Here $x\geq5/6$ subsumes the root-triangle constraint $x\geq2/3$.

The lower boundary of \eqref{eq:split-dual-region} has vertices
\[
  (2,0),\qquad \left(\frac43,\frac13\right),\qquad
  \left(\frac56,\frac56\right).
\]
The dual objective is $\binom{k}{2}x+k\ell y$, so its minimum is attained at
one of these vertices.  Subtracting the three possible values from
$e(S_{k,\ell})=\binom{k}{2}+k\ell$ gives
\eqref{eq:split-functional}.
\end{proof}

The three expressions in \eqref{eq:split-functional}, maximized over the real
variable $k$, are bounded respectively by
\begin{equation}\label{eq:three-global-bounds}
  Q(n),\qquad \frac{(4n+1)^2}{120},qquad
  \frac{n(n-1)}{12}.
\end{equation}
The first bound is the square identity \eqref{eq:square-identity}; the other
two follow by completing the square and by $e(S_{k,\ell})\leq\binom n2$.
If $\ell=0$ and $k\geq4$, then
$\PhiF(K_k)=e(K_k)/6$: the constant dual weight $5/6$ is optimal, as witnessed
by the uniform fractional packing of all copies of $K_4$.  If $k\leq3$, then
$\PhiF(S_{k,n-k})\leq e(S_{k,n-k})\leq3n$.

For $n\geq100$, every branch except the first in
\eqref{eq:split-functional}, including the preceding degenerate cases, is
strictly smaller than
\begin{equation}\label{eq:branch-gap}
  \frac{n^2}{6}-\frac{n^2}{40}.
\end{equation}
For the second branch this is equivalent to $n^2>8n+1$, while
$3n<n^2/6-n^2/40$ follows from $17n>360$; the third branch and the complete
graph are smaller still.

\begin{corollary}\label{cor:global-Phi-bound}
Every chordal graph $G$ of order $n\geq100$ satisfies
\[
  \PhiF(G)\leq Q(n).
\]
\end{corollary}

\begin{proof}
Apply \cref{prop:fine-symmetrization}.  The functional does not decrease along
the path, and \eqref{eq:three-global-bounds} bounds its complete-split
endpoint by $Q(n)$.
\end{proof}

We shall also need a quantitative statement about a nearly extremal terminal
graph.

\begin{lemma}\label{lem:terminal-stability}
Let $n\geq100$, and suppose that a complete-split graph $S$ of order $n$
satisfies
\[
  \PhiF(S)\geq\frac{n^2}{6}-\eta n^2.
\]
If $\mathcal{D}=Q(n)-\PhiF(S)$, then the first branch in
\eqref{eq:split-functional} gives the value of $\PhiF(S)$ and
\begin{equation}\label{eq:terminal-stability}
  \distsp(S)\leq n\sqrt{\frac{2\mathcal{D}}3}+n.
\end{equation}
\end{lemma}

\begin{proof}
Since $\eta=10^{-30}<1/40$, the branch separation
\eqref{eq:branch-gap} forces the first branch.  If its root has order $k$,
the square identity \eqref{eq:square-identity} gives
\[
  \left|k-\frac{2n+1}{6}\right|
  \leq\sqrt{\frac{2\mathcal{D}}3}.
\]
Changing the root order to $\lfloor n/3\rfloor$ costs at most $n$ edits per
vertex-role change.  The offset between $(2n+1)/6$ and $\lfloor n/3\rfloor$
is less than one, proving \eqref{eq:terminal-stability}.
\end{proof}

\section{Global stability by first entry}

\begin{theorem}[fixed-neighbourhood stability]\label{thm:global-stability}
Let $G$ be chordal of order $n\geq10^{32}$.  If
\begin{equation}\label{eq:near-extremal-Phi}
  \PhiF(G)\geq\frac{n^2}{6}-\eta n^2,
\end{equation}
then
\begin{equation}\label{eq:global-stability-conclusion}
  \distsp(G)<\frac{\rho n^2}{2}.
\end{equation}
\end{theorem}

\begin{proof}
Take the path $G=G_0,\ldots,G_L$ from
\cref{prop:fine-symmetrization}.  Its functional values are no smaller than
$\PhiF(G)$.  Therefore, for every graph $H$ on the path,
\begin{align}
  Q(n)-\PhiF(H)
  &\leq \eta n^2+\frac n6+\frac1{24}
   \leq2\eta n^2,
  \label{eq:path-deficit}
\end{align}
where the last inequality follows from $n\geq10^{32}$.
By \cref{lem:terminal-stability}, the endpoint satisfies
\[
  \frac{\distsp(G_L)}{n^2}
  \leq\sqrt{\frac{4\eta}{3}}+\frac1n
  <\frac\rho4.
\]

Suppose, contrary to \eqref{eq:global-stability-conclusion}, that
$\distsp(G_0)\geq\rho n^2/2$.  Let $i$ be the first path index with
$\distsp(G_i)<\rho n^2/2$.  A path step changes at most $n-2$ edges, and hence
\begin{equation}\label{eq:first-entry-lower}
  \frac{\rho n^2}{2}-n
  \leq\distsp(G_i)<\frac{\rho n^2}{2}.
\end{equation}
In particular, $G_i$ lies inside the radius-$\rho n^2$ neighbourhood, so
\cref{lem:local-stability} applies.  Combining
\eqref{eq:local-stability} with \eqref{eq:path-deficit} gives
\begin{align}
  \frac{\distsp(G_i)}{n^2}
  &\leq18\eta+\sqrt{\frac{4\eta}{3}}+\frac1n \notag\\
  &<18\cdot10^{-30}+2\cdot10^{-15}+10^{-32}
   <\frac\rho4.
  \label{eq:first-entry-upper}
\end{align}
But $\rho/2-1/n>\rho/4$, so
\eqref{eq:first-entry-lower} and \eqref{eq:first-entry-upper} are
incompatible.  This proves the theorem.
\end{proof}

\section{Proof of the main theorem}

We now combine the near and far regimes.  Recall the definition of $N$ in
\eqref{eq:N-definition}.

\begin{theorem}\label{thm:eventual-upper}
Every chordal graph $G$ on $n\geq N$ vertices satisfies
\[
  \cple{4}(G)\leq M(n).
\]
\end{theorem}

\begin{proof}
If \eqref{eq:near-extremal-Phi} holds, then
\cref{thm:global-stability} places $G$ within the radius-$\rho n^2$
neighbourhood.  Equation \eqref{eq:local-integral-conclusion} gives the
stronger bound $\cple{3}(G)\leq M(n)$.

Otherwise, apply \cref{thm:packing-transfer} with $\varepsilon=\eta/2$.
Using \eqref{eq:cp4-identity},
\begin{align*}
  \cple{4}(G)
  &=\PhiF(G)+\Wfrac(G)-\Wint(G)\\
  &<\frac{n^2}{6}-\eta n^2+\frac{\eta n^2}{2}
   <\frac{n^2}{6}\leq M(n),
\end{align*}
where the final inequality is valid for the present range of $n$.  This
completes both cases.
\end{proof}

It remains to match the upper bound.  We record the standard complete-split
calculation because it also explains the linear term in $M(n)$.

\begin{lemma}\label{lem:split-exact-cp}
If $\ell\geq\chi'(K_k)$, then
\begin{equation}\label{eq:split-exact-cp}
  \cp(S_{k,\ell})=k\ell-\binom k2.
\end{equation}
\end{lemma}

\begin{proof}
Properly edge-colour the clique $K_k$ using at most $\ell$ colours, identified
with distinct vertices of the independent side.  For every root edge $xy$,
use the triangle consisting of $x,y$ and the vertex corresponding to the
colour of $xy$.  Properness makes these triangles edge-disjoint.  Cover every
unused spoke by a $K_2$.  The resulting partition has
\[
  \binom k2+\left(k\ell-2\binom k2\right)
  =k\ell-\binom k2
\]
parts.

For the reverse inequality, assign weight $-1$ to each root edge and weight
$+1$ to each spoke.  Every clique in a complete-split graph has total edge
weight at most $1$: a clique contains at most one independent-side vertex,
and, if it contains $r$ root vertices and one outside vertex, its weight is
$r-\binom r2\leq1$.  Its value is nonpositive when it contains only root
vertices.  Summing over any clique partition shows that the number of parts is
at least the total edge weight $k\ell-\binom{k}{2}$.
\end{proof}

Choose an integer $k$ nearest to $(2n+1)/6$ and put $\ell=n-k$.  Then
$\ell\geq k\geq\chi'(K_k)$, using Vizing's theorem, and a check of the three
residue classes of $n$ modulo $3$ gives
\begin{equation}\label{eq:integer-maximization}
  k(n-k)-\binom k2
  =\left\lfloor\frac{n(n+1)}6\right\rfloor=M(n).
\end{equation}
Together with \cref{thm:eventual-upper}, this proves
\cref{thm:main}.

For completeness, let $C=N$.  If $0<n<N$, the partition into individual
edges gives
\[
  \cp(G)\leq\binom n2\leq\frac{n^2}{6}+Nn.
\]
For $n\geq N$, \cref{thm:eventual-upper} and
$M(n)\leq n^2/6+n/6$ give the same conclusion.  The empty graph is immediate.
This proves \cref{cor:erdos81}.

\section{Verification and scope}

The proof uses exactly three non-elementary external inputs:
\cref{thm:vizing,thm:HJ,thm:packing-transfer}.  In particular, it does not
assume a linear bound for the general integral--fractional packing gap.  The
transfer theorem is invoked only in the far regime, where a fixed quadratic
deficit pays for its $o(n^2)$ loss.  In the near regime,
\cref{lem:terminal-construction} constructs an actual edge partition into
triangles and $K_2$'s.

The accompanying repository contains exact rational checks of the numerical
inequalities, finite mixed-packing certificates, and the single-copy
mechanism.  These computations are regression tests; none is a premise of the
proof for arbitrary $n$.  The repository also records the precise boundary of
the staged Lean formalization and the commands used to reproduce every
machine check.

\subsection*{AI-assisted development disclosure}

AI systems were used in exploration, proof organization, adversarial review,
finite computation, and editorial work.  The mathematical claims have been
restated as the self-contained argument above, with all external results
identified explicitly.  Responsibility for verification, attribution, and
publication remains with the human authors; no AI system is an author.

\bibliographystyle{alpha}
\bibliography{references}

\end{document}
